\section*{Author Section 3.1: dimension of vanishing spaces}
In this section, we will be interested in studying the vanishing spaces and closures of a random set $Z$ obtained by picking $K=(1-\epsilon)\binom m{\le r}$ points from $\mathbb F_2^m$ (where $\epsilon>0$ is some constant). We will use the following lemma to prove the above theorem.
\paragraph{Lemma 3.4 (size of degree $r$-vanishing space of random set).} There exists a constant $\gamma_0>0$ such that for all $\epsilon>0$ the following is true. Let $K=(1-\epsilon)\binom m{\le r}$ and $r<\gamma_0m$. Then,
\[\mathbb E\left[\frac{|I_r(Z)|}{\operatorname{exp}_2\left(\binom m{\le r}\right)}\right]=2^{-K}(1+o(1))\]
where $Z=\{z_1,\ldots,z_K\}$ is a uniformly random set of $K$ distinct points in $\mathbb F_2^m$. We first use the above lemma to prove Theorem 1.1. Lemma 3.4 is proved in Subsection 3.2.
\paragraph{Proof of Theorem 1.1.} Let $p$ be the probability that $\dim(I_r(Z))=\binom m{\le r}-K$. Note that $\dim(I_r(Z))$ is always at least $\binom m{\le r}-K$, as the space $I_r(Z)$ is a subspace of $\mathrm{RM}(m,r)$ defined by $K$ linear equations. Thus,
\begin{align*}\mathbb E\left[\frac{|I_r(Z)|}{\operatorname{exp}_2\left(\binom m{\le r}\right)}\right]&=\mathbb E\left[\operatorname{exp}_2\left(\dim(I_r(Z))-\binom m{\le r}\right)\right]\\&\ge p\operatorname{exp}_2(-K)+(1-p)2\operatorname{exp}_2(-K).\end{align*}
Combining this with Lemma 3.4 we get that
\[p\operatorname{exp}_2(-K)+(1-p)2\operatorname{exp}_2(-K)\le2^{-K}(1+o(1))\]
which gives $1-p\le o(1)$. Hence, we get that $p=1-o(1)$.\qed
\section*{Author Section 3.2: proof of Lemma 3.4}
We start with reformulating the problem of bounding size of $I_r(Z)$ to computing a certain weighted sum of the polynomials in $\mathrm{RM}(m,r)$. Let $Q$ be a uniformly random polynomial chosen from $\mathrm{RM}(m,r)$. Then,
\begin{align*}
\mathbb E\left[\frac{|I_r(Z)|}{\operatorname{exp}_2\left(\binom m{\le r}\right)}\right]&=\mathbb E_{Z,Q}[\mathbf1[Q\in I_r(Z)]]\\&=\mathbb E_Q[\mathbb E_Z[\mathbf1[Q\in I_r(Z)]]]\\
&=\sum_{P\in\mathrm{RM}(m,r)}\operatorname{exp}_2\left(-\binom m{\le r}\right)\frac{\binom{(1-\mathrm{wt}(P))2^m}K}{\binom{2^m}K}\\
&=\operatorname{exp}_2(-K)\left[\sum_{P\in\mathrm{RM}(m,r)}\operatorname{exp}_2\left(-\epsilon\binom m{\le r}\right)\frac{\binom{(1-\mathrm{wt}(P))2^m}K}{\binom{2^m}K}\right].
\end{align*}
To complete the proof of Lemma 3.4 we therefore need the following claim.
\paragraph{Claim 3.5.}
\[\sum_{P\in\mathrm{RM}(m,r)}\operatorname{exp}_2\left(-\epsilon\binom m{\le r}\right)\frac{\binom{(1-\mathrm{wt}(P))2^m}K}{\binom{2^m}K}\le1+o(1).\]
Given Claim 3.5, we see that
\begin{align*}\mathbb E\left[\frac{|I_r(Z)|}{\operatorname{exp}_2\left(\binom m{\le r}\right)}\right]&=\operatorname{exp}_2(-K)\left[\sum_{P\in\mathrm{RM}(m,r)}\operatorname{exp}_2\left(-\epsilon\binom m{\le r}\right)\frac{\binom{(1-\mathrm{wt}(P))2^m}K}{\binom{2^m}K}\right]\\&\le\operatorname{exp}_2(-K)(1+o(1)),\end{align*}
which concludes the proof of Lemma 3.4. Next, we proceed to prove Claim 3.5.
\paragraph{Proof of Claim 3.5.} Let $u=\binom m{\le r}$ in the following. Recall that $K=(1-\epsilon)u$. Also,
\[\sum_{P\in\mathrm{RM}(m,r)}\operatorname{exp}_2(-\epsilon u)\frac{\binom{(1-\mathrm{wt}(P))2^m}K}{\binom{2^m}K}\le\sum_{P\in\mathrm{RM}(m,r)}\operatorname{exp}_2(-\epsilon u)(1-\mathrm{wt}(P))^{(1-\epsilon)u}.\]
Set $\delta=\left(1/\binom m{\le r}\right)^2$ and let $\mathrm{RM}_\delta(m,r)=\{P\in\mathrm{RM}(m,r):|\mathrm{wt}(P)-1/2|\ge\delta/2\}$, i.e., the set of polynomials with bias at least $\delta$. Now,
\begin{multline*}\sum_{P\in\mathrm{RM}_\delta(m,r)}\operatorname{exp}_2(-\epsilon u)(1-\mathrm{wt}(P))^{(1-\epsilon)u}\\\le2\sum_{P:\mathrm{wt}(P)\le1/2-\delta/2}\operatorname{exp}_2(-\epsilon u)(1-\mathrm{wt}(P))^{(1-\epsilon)u},\end{multline*}
since the term corresponding to a polynomial $P$ of weight greater than $1/2+\delta/2$ can be upper bounded by the term corresponding to the polynomial $1+P$ which has weight at most $1/2-\delta/2$.
Since the RHS of the above equation involves polynomials whose weights are in the interval $[1/2^r,1/2-\delta]$, we will split this interval into sub-intervals and analyse the contribution from each.
\begin{align*}\mathrm{Low}_i&:=[1/2^{i+1},1/2^i]&&\text{for }i=2,3,\ldots,r-1,\\\mathrm{Med}_i&:=[(1/2-1/2^i),(1/2-1/2^{i+1})]&&\text{for }i=2,3,\ldots,t=\log(1/\delta).\end{align*}
Let us use the following quantities to denote the number of polynomials with weights in the above intervals:
\[L_i:=|\{P\in\mathrm{RM}(m,r):\mathrm{wt}(P)\in\mathrm{Low}_i\}|,\qquad M_i:=|\{P\in\mathrm{RM}(m,r):\mathrm{wt}(P)\in\mathrm{Med}_i\}|.\]
Therefore,
\begin{align*}&\sum_{P\in\mathrm{RM}_\delta(m,r)}\operatorname{exp}_2(-\epsilon u)(1-\mathrm{wt}(P))^{(1-\epsilon)u}\\&\quad\le2\sum_{P:\mathrm{wt}(P)\le1/2-\delta/2}\operatorname{exp}_2(-\epsilon u)(1-\mathrm{wt}(P))^{(1-\epsilon)u}\\&\quad\le2\left(\sum_{i=2}^{r-1}L_i(1-1/2^{i+1})^k+\sum_{i=2}^tM_i(1/2+1/2^i)^k\right)\frac1{\operatorname{exp}_2\left(\binom m{\le r}-K\right)}.\end{align*}
By Claim 3.6 proved below using the weight distribution bounds of Sberlo and Shpilka (Theorem 2.1 and Theorem 2.2), we have that
\[\left(\sum_{i=2}^{r-1}L_i(1-1/2^{i+1})^k+\sum_{i=2}^tM_i(1/2+1/2^i)^k\right)\frac1{\operatorname{exp}_2\left(\binom m{\le r}-K\right)}\le O(\delta^2).\]
Hence,
\begin{align*}&\sum_{P\in\mathrm{RM}(m,r)}\operatorname{exp}_2(-\epsilon u)(1-\mathrm{wt}(P))^{(1-\epsilon)u}\\&=\sum_{P\in\mathrm{RM}_\delta(m,r)}\operatorname{exp}_2(-\epsilon u)(1-\mathrm{wt}(P))^{(1-\epsilon)u}\\&\qquad+\sum_{P\notin\mathrm{RM}_\delta(m,r)}\operatorname{exp}_2(-\epsilon u)(1-\mathrm{wt}(P))^{(1-\epsilon)u}\\&\le O(\delta^2)+\sum_{P\notin\mathrm{RM}_\delta(m,r)}\operatorname{exp}_2(-\epsilon u)(1/2+\delta/2)^{(1-\epsilon)u}\\&\le O(\delta^2)+(1+\delta)^{(1-\epsilon)u}\\&\le O(\delta^2)+\exp(\delta u)\\&\le1+O(1/u)=1+o(1).\end{align*}\qed
It remains to prove the following technical claim.
\paragraph{Claim 3.6.} 1. For all $i=2,\ldots,r-1$, we have
\[\frac{L_i(1-1/2^{i+1})^k}{\operatorname{exp}_2\left(\binom m{\le r}-k\right)}\le\delta^3.\]
2. For all $i=2,\ldots,t=\log(1/\delta)$, we have
\[\frac{M_i(1/2+1/2^i)^k}{\operatorname{exp}_2\left(\binom m{\le r}-k\right)}\le\delta^3.\]
\paragraph{Prooff of Claim 3.6(1).} Note that $(1-1/2^{i+1})^k\le\exp(-k/2^{i+1})\le\operatorname{exp}_2(-k/2^{i+1})$. Hence,
\[\frac{L_i(1-1/2^{i+1})^k}{\operatorname{exp}_2\left(\binom m{\le r}-k\right)}\le\frac{\mathrm{WtDist}_{m,r}(2^{-i})\operatorname{exp}_2(-k/2^{i+1})}{\operatorname{exp}_2\left(\binom m{\le r}-k\right)}.\]
Using Theorem 2.1 to bound $\mathrm{WtDist}_{m,r}(2^{-i})$, we get
\begin{align*}\frac{L_i(1-1/2^{i+1})^k}{\operatorname{exp}_2\left(\binom m{\le r}-k\right)}&\le\operatorname{exp}_2\left(O(m^4)+O(i\gamma^{i-1})\binom m{\le r}-k/2^{i+1}-\binom m{\le r}+k\right)\\&=\operatorname{exp}_2\left(O(m^4)-\binom m{\le r}(1-O(i\gamma^{i-1}))+k(1-2^{-i-1})\right),\end{align*}
where $\gamma=r/m$. If this $\gamma$ is a small enough absolute constant, we can see that the $O(i\gamma^{i-1})$ term is always at most $2^{-i-1}$. Using this and continuing the computation above, we get
\begin{align*}\frac{L_i(1-1/2^{i+1})^k}{\operatorname{exp}_2\left(\binom m{\le r}-k\right)}&\le\operatorname{exp}_2\left(O(m^4)-\left(\binom m{\le r}-k\right)(1-2^{-i-1})\right)\\&\le\operatorname{exp}_2\left(O(m^4)-\epsilon\binom m{\le r}(1-2^{-i-1})\right)\\&\le\operatorname{exp}_2\left(O(m^4)-\frac\epsilon2\binom m{\le r}\right)\\&\le\operatorname{exp}_2\left(-\Omega\left(\binom m{\le r}\right)\right)\le\delta^3\end{align*}
where for the second inequality above, we used the fact that $k\le K=(1-\epsilon)\binom m{\le r}$.\qed
\paragraph{Proof of Claim 3.6(2).} Note that
\begin{align*}(1/2+1/2^i)^k&=2^{-k}(1+1/2^{i-1})^k\\&\le2^{-k}\exp(k/2^{i-1})=2^{-k}\operatorname{exp}_2(k/2^{i-1}\cdot\log_2e)\\&=\operatorname{exp}_2(-k(1-2\log_2e\cdot2^{-i})).\end{align*}
Using the fact that $M_i\le\mathrm{WtDist}_{m,r}(1/2-1/2^{i+1})$, we get
\begin{align*}&\frac{\mathrm{WtDist}_{m,r}(1/2-1/2^{i+1})\operatorname{exp}_2(-k(1-2\log_2e\cdot2^{-i}))}{\operatorname{exp}_2\left(\binom m{\le r}-k\right)}\\&\quad\le\frac{\operatorname{exp}_2\left(O(m^4)+(1-2^{-O(\max\{\gamma^2i,\gamma\})})\binom m{\le r}-k(1-2\log_2e\cdot2^{-i}))\right)}{\operatorname{exp}_2\left(\binom m{\le r}-k\right)}\\&\quad=\operatorname{exp}_2\left(O(m^4)-\frac1{2^{O(\max\{\gamma^2i,\gamma\})}}\binom m{\le r}+k\cdot\frac1{2^{-(i-1-\log_2\log_2e)}}\right)\end{align*}
where $\gamma=r/m$, and we used Theorem 2.2 for the second inequality.
Note that as long as $\gamma$ is a small enough absolute constant, the $O(\max\{\gamma^2i,\gamma\})$ term above is at most $i-1-\log_2\log_2e$ for each $i\in\{2,\ldots,t\}$. Using this bound, we continue the above computation as follows.
\begin{align*}\frac{M_i(1/2+1/2^i)^k}{\operatorname{exp}_2\left(\binom m{\le r}-k\right)}&\le\operatorname{exp}_2\left(O(m^4)-2^{-O(\max\{\gamma^2i,\gamma\})}\left(\binom m{\le r}-k\right)\right)\\&\le\operatorname{exp}_2\left(O(m^4)-2^{-O(\max\{\gamma^2i,\gamma\})}\left(\binom m{\le r}-K\right)\right)\\&=\operatorname{exp}_2\left(O(m^4)-2^{-O(\max\{\gamma^2i,\gamma\})}\epsilon\binom m{\le r}\right)\\&\le\operatorname{exp}_2\left(O(m^4)-\Omega\left(\sqrt{\binom m{\le r}}\right)\right)\le\delta^3,\end{align*}
where for the last inequality we used the fact that $\binom m{\le r}\ge2^{\Omega(\gamma\log(1/\gamma)m)}$ and hence for small enough absolute constant $\gamma$, we have
\begin{align*}2^{-O(\max\{\gamma^2i,\gamma\})}\binom m{\le r}&\ge2^{-O(\gamma^2t)}\binom m{\le r}\\&\ge\operatorname{exp}_2(\Omega(\gamma\log(1/\gamma)m)-O(\gamma^2\log(1/\delta)))\\&=\operatorname{exp}_2(\Omega(\gamma\log(1/\gamma)m)-O(\gamma^3\log(1/\gamma)m))\\&=\operatorname{exp}_2(\Omega(m)).\end{align*}\qed
This completes the proof of Lemma 3.4.
